\documentclass[10pt,a4paper]{amsart}
\usepackage[margin=19mm]{geometry}
\usepackage{amsmath,amssymb,mathtools}
\usepackage[scr=rsfs,cal=euler]{mathalfa}
\usepackage{xfrac}
\usepackage[unicode,hidelinks,bookmarks=false]{hyperref}
\newcommand{\ie}{i.e.\ }
\newcommand{\cf}{cf.\ }
\newcommand{\resp}{respectively\ }
\newcommand{\Subsection}{\subsection}
\providecommand{\nobreakdash}{\nobreak\hbox{-}}
\setlength{\emergencystretch}{2em}
\multlinegap0pt
\usepackage{tikz}
\usepackage{amssymb}
\usepackage{xcolor}
\usepackage[shortlabels]{enumitem}
\usepackage{float}
\newtheorem{proposition}{Proposition}
\newtheorem{lemma}{Lemma}
\newcommand{\F}{\mathbb F}
\title{On Erd\H{o}s--Falconer distance problem in even dimensions}
\author{Thang Pham, Chun-Yen Shen, Boqing Xue}
\date{Source: arXiv:2607.17324v1 (CC BY 4.0)}
\begin{document}
\maketitle

\noindent\emph{Source and licence.} On Erd\H{o}s--Falconer distance problem in even dimensions. Version arXiv:2607.17324v1 (CC BY 4.0), distributed by arXiv under the Creative Commons Attribution 4.0 International licence, http://creativecommons.org/licenses/by/4.0/, as recorded in the arXiv licence metadata for this exact version and shown on the abstract page https://arxiv.org/abs/2607.17324v1. Full licence notice: \texttt{licenses/arxiv-2607.17324-CC-BY-4.0.txt}.\par\smallskip

\noindent\textit{Editorial scope.} Counted unit: the article's proposition prop:prime-planar-both-2 (source0.tex lines 780--787). The article's complete original proof of it is delivered with it (source0.tex lines 1223--1299, one \texttt{proof} environment of 77 lines, which the article places away from the statement). No mathematics is written, repaired or re-derived: the statement and the proof are the article's own lines, in the article's own notation, and no retained line is substituted or re-flowed.  Retained uncounted as the source-local context the proof needs: lines 820--826; lines 867--874.  Nothing was omitted from any retained span, so the package carries no omission record.  No retained line contains a citation, so the wrapper carries no bibliography and no bibliography entry.  No retained line contains a figure environment, an image-inclusion command or drawing code, so no image is delivered and the package declares no asset dependency.  Editorial formatting only, in addition: an AMS \texttt{amsart} preamble that restates the article's own declarations and macros used by the retained lines, namely the environments \texttt{proposition}, \texttt{lemma} on separate counters, together with the standard packages those lines need.  The retained environments print in this document as Proposition 1, Lemma 1 in order of appearance, whereas the article prints the same items with the numbering of its own full document.  The complete native source of the article as distributed in the arXiv e-print for this exact version is bundled unchanged as \texttt{sources/205557/source0.tex}.
\begin{lemma}
\label{lem:split-fibre-pruning}
Let $A\subset\F_p^2$, $|A|\ge2p$, and $0<\delta<1$.  If $|\Delta_{H,a}(A)|<\delta p$ for every $a\in A$, then
\[
M_0(A)\le \frac{2\delta}{1-\delta}\,\frac{p^2}{|A|}.
\]
\end{lemma}

\begin{lemma}
\label{lem:split-isosceles}
There is an absolute constant $\widetilde{C}>0$ such that, if
$p\le |A|\le p^{4/3}$, then
\begin{equation}\label{eq:split-isosceles}
\mathcal T_H(A) \le \frac{|A|^3}{p}+\widetilde{C}\left(p^{2/3}|A|^{5/3}+p^{1/4}|A|^2 + p^{1/2}M_0(A)^{1/2}|A|^{3/2}\right).
\end{equation}
\end{lemma}

\begin{proposition}
\label{prop:prime-planar-both-2}
There are absolute constants $C_4,c_4>0$ such that, for every odd
prime $p$, and every $A\subset\F_p^2$,
\begin{equation} \label{eq:prime-planar-both}
|A|\ge C_4 p^{5/4} \quad\Longrightarrow\quad \Delta_{\mathrm{pin},H}(A)\ge c_4 p.
\end{equation}
\end{proposition}

\begin{proof}[Proof of Proposition \ref{prop:prime-planar-both-2}]
Enlarge the eventual choice of $C_4$, if necessary, so that
$C_4\ge2$; hence $|A|\ge C_4p^{5/4}$ implies $|A|\ge2p$, as required
in Lemma~\ref{lem:split-fibre-pruning}.  It suffices, for all
sufficiently large $p$, to consider
\[
 C_4p^{5/4}\le |A|\le p^{4/3},
\]
because, once $C_4$ is fixed, a larger $A$ has a subset of size
$\lfloor p^{4/3}\rfloor\ge C_4p^{5/4}$.  A pin found in that subset
is also a pin of the original set, and its pinned distance set can
only increase.  We include the displayed size inequality in the
large-prime cutoff and treat the remaining primes at the end.
Fix a sufficiently small absolute
$0<\delta<1/16$ and suppose, for a contradiction, that every pin
determines fewer than $\delta p$ values.  Lemma~\ref{lem:split-fibre-pruning}
gives
\[
 M_0(A)\le
 \frac{2\delta}{1-\delta}\frac{p^2}{|A|}.
\]
Lemma~\ref{lem:split-isosceles} therefore yields
\begin{equation}
 \label{eq:TH-upper-bad}
 \mathcal T_H(A)
 \le\frac{|A|^3}{p}
 +\widetilde{C}\left(
 p^{2/3}|A|^{5/3}
 +p^{1/4}|A|^2
 +\sqrt{\frac{\delta}{1-\delta}}\,p^{3/2}|A|
 \right)
\end{equation}
with an absolute constant $\widetilde{C}$.

%For $a\in A$, put $\nu_a(t)=|\{b\in A:H(a-b)=t\}|$. 
The zero cone is the union of a horizontal and a vertical line, so
$\nu_a(0)\le2p-1$.  If
$H(a-b)=H(a-c)\ne0$ and $H(b-c)=0$, then $b=c$:
for example, if $b-c=(d,0)$ with $d\ne0$, subtracting the two
equalities gives $d(a_2-c_2)=0$, which forces the common value to be
zero.  Thus
\[
 \mathcal T_H(A)+|A|^2
 \ge\sum_{a\in A}\sum_{t\ne0}\nu_a(t)^2.
\]
Applying the Cauchy-Schwarz inequality for each fixed \(a\in A\),
then summing over \(a\in A\), the bad-pin hypothesis and the preceding
estimate give
\begin{equation}
 \label{eq:TH-lower-bad}
 \mathcal T_H(A)
 \ge
 \frac{|A|(|A|-2p)^2}{\delta p}-|A|^2.
\end{equation}

At $|A|\ge C_4p^{5/4}$, the three error terms in \eqref{eq:TH-upper-bad}, divided by $|A|^3/p$, are respectively
\[
 O(C_4^{-4/3}),\qquad
 O(C_4^{-1}),\qquad
 O\!\left(\sqrt{\delta}\,C_4^{-2}\right).
\]
After fixing $0<\delta<1/16$, choose $C_4$ sufficiently large so
that the sum of the three normalized error
terms above is at most $1$.  Then
\eqref{eq:TH-upper-bad} is at most $2|A|^3/p$.  On the other hand, for
all sufficiently large $p$ the inequality
$|A|\ge C_4p^{5/4}$ gives $|A|-2p\ge |A|/2$, and
\eqref{eq:TH-lower-bad} yields
\[
 \frac{\mathcal T_H(A)}{|A|^3/p}
 \ge \frac1{4\delta}-\frac{p}{|A|}>2.
\]
This is the required contradiction. 
Let $p_0$ exceed both this cutoff and the subset-reduction
cutoff above.  For $p<p_0$, every nonempty pinned distance set contains
$0$ and hence has size at least $1\ge p/p_0$. Taking $C_4=\min\{\delta,p_0^{-1}\}$ proves the assertion for every odd prime.
\end{proof}

\end{document}
